Scaling the four named cases to the whole missed population needs one methodological care: a SHAP value and a distance-to-threshold are commensurable only in margin (log-odds) space, where TreeSHAP's additivity holds, and a comparison run in probability space would be a category error with plausible-looking output. Taken there (Table~\ref{tab:vndeficit}), the two natural readings of ``the prior caused the miss'' bracket the answer from opposite sides. The additive accounting is the conservative one: on $24$ of the $141$ missed \texttt{.vn} rows ($0.170$, Wilson 95\% CI $0.117$--$0.241$) the \texttt{tld\_len} contribution alone is at least as large as the row's entire deficit (median deficit $1.22$ against a median contribution of $0.89$). The intervention is the generous one: re-scoring those rows with \texttt{tld\_len} pinned to the fit window's non-\texttt{.vn} mode ($3$) and nothing else touched carries $121$ of $141$ ($0.858$, $0.791$--$0.906$) across the threshold, a median gain of $+1.81$ in margin. The direction of that gap is itself the finding, and it is the one Section~\ref{ssec:suffixprior} predicts: Shapley splits an interaction between its partners, so a prior carried \emph{jointly} (which is exactly what a two-character suffix length must be, since it cannot identify \texttt{.vn} on its own) is under-credited by the decomposition and recovered in full by the intervention (Spearman $0.96$ between the two per row). Neither number is the whole causal story: the intervention moves one feature while \texttt{url\_len} and \texttt{dom\_len} still count the original suffix's characters, and the crossing count does depend on the value chosen ($129$ at \texttt{tld\_len}$=4$, $82$ at $5$). What survives both readings is that the misses are a suffix-prior phenomenon rather than an absence of phishing signal in the URL%
